\begingroup
\setlength{\parindent}{0pt}
\setlength{\parskip}{0.40\baselineskip}
\setlength{\abovedisplayskip}{6.5pt plus 1.5pt minus 1.5pt}
\setlength{\belowdisplayskip}{6.5pt plus 1.5pt minus 1.5pt}
\setlength{\jot}{4.5pt}

\section*{How the Proof Works}

For the global three-dimensional Navier--Stokes regularity problem, take one
smooth finite-energy flow with fixed viscosity \(\nu>0\):
\[
 \partial_tu+(u\!\cdot\nabla)u+\nabla p=\nu\Delta u,
 \qquad \nabla\!\cdot u=0,
 \qquad u(0)=u_0.
\]
The danger is that this flow could sharpen one of its own structures beyond
smooth continuation in finite time.  A vortex can concentrate velocity
gradients while it stretches.
As surrounding strain pulls the tube along its axis, incompressibility forces
it to contract across its width: a longer piece of fluid cannot keep its old
cross-section without gaining volume.  Aligned stretching can also strengthen
the rotation.  The core narrows, the velocity turns more sharply across it,
and the gradients rise.  This nonlinear sharpening is opposed by viscosity,
which spreads the concentrated vorticity into its surroundings.  Both effects
can grow stronger as the core gets smaller.

To shrink a width by a factor \(\lambda>1\) at the Navier--Stokes scale,
multiply velocity by \(\lambda\) and shorten time by \(\lambda^2\):
\[
 u_\lambda(x,t)=\lambda u(\lambda x,\lambda^2t),
 \qquad
 p_\lambda(x,t)=\lambda^2p(\lambda x,\lambda^2t).
\]
Nonlinear transport and viscosity both gain a factor \(\lambda^3\): viscosity
has not gained a scaling advantage over the sharpening.  Equal scaling weight
does not tell us which effect wins dynamically, or give a bound on either one.

The faster motion occupies less volume.  Velocity squared gains
\(\lambda^2\), while volume loses \(\lambda^3\), so kinetic energy scales down
by \(\lambda^{-1}\).  It can remain finite as velocity and gradients grow.
The Sobolev norms register this concentration at different spatial heights:
\[
 \|u_\lambda(\cdot,t)\|_{\dot H^s}
 =\lambda^{s-1/2}
  \|u(\cdot,\lambda^2t)\|_{\dot H^s}.
\]
Below \(s=\tfrac12\), the scaled norm falls; above it, the norm rises.  At
\(s=\tfrac12\), its size is unchanged by this concentration.  That is the
critical height: \(\dot H^{1/2}\) is scaling-critical, not conserved in time.

If the vortex kept repeating this narrowing and acceleration, its width,
velocity scale, and duration at stage \(j\) would follow
\[
\begin{gathered}
 \ell_j=\lambda^{-j}\ell_0,
 \qquad
 U_j=\lambda^jU_0,
 \qquad
 \tau_j=\lambda^{-2j}\tau_0,
 \\
 \sum_{j=0}^{\infty}\tau_j
 =\frac{\tau_0}{1-\lambda^{-2}}<\infty.
\end{gathered}
\]
The widths collapse and the velocity scale diverges, yet infinitely many
stages fit into a finite interval.  This is the hypothetical Zeno cascade:
finite energy can survive while the core becomes arbitrarily narrow and fast.

In this picture, in order for a singularity to happen, velocity
has to blow up.  In order for that to happen, its acceleration would also have
to blow up, which means its jerk has to blow up, and so on, and so on:
\[
 u,\qquad \partial_tu,\qquad \partial_t^2u,\qquad\ldots
\]
The hypothetical cascade cannot stop at any fixed temporal order.

The equation carries this repeated temporal climb into its nonlinear and
viscous structure:
\[
\begin{aligned}
 \partial_tu
 &=\nu\Delta u-(u\!\cdot\nabla)u-\nabla p,\\
 \partial_t^2u
 &=\nu\Delta\partial_tu
   -(\partial_tu\!\cdot\nabla)u
   -(u\!\cdot\nabla)\partial_tu
   -\nabla\partial_tp,\\
 \partial_t^3u
 &=\nu\Delta\partial_t^2u
   -(\partial_t^2u\!\cdot\nabla)u
   -2(\partial_tu\!\cdot\nabla)\partial_tu
   -(u\!\cdot\nabla)\partial_t^2u
   -\nabla\partial_t^2p.
\end{aligned}
\]
The Laplacian follows the rising temporal response, while the nonlinear term
opens into every product allowed by the product rule.  After three lines the
bookkeeping is already hiding the relation, so put
\[
 V_r:=\partial_t^ru,
 \qquad
 Q_r:=\partial_t^rp.
\]
At order \(r\), those product-rule copies occur with multiplicity
\(\binom ra\), so the expansion closes as
\[
\begin{aligned}
 V_{r+1}
 &+\sum_{a=0}^{r}\binom ra
   (V_a\!\cdot\nabla)V_{r-a}
   +\nabla Q_r
   =\nu\Delta V_r,\\
 Q_r
 &=R_iR_j\sum_{a=0}^{r}\binom ra
   V_{a,i}V_{r-a,j}.
\end{aligned}
\]
Incompressibility fixes \(Q_r\) through the Riesz transforms \(R_i\), so the
pressure row stays with the nonlinear row at every order.

If a singularity were forming, one would expect spatial regularity to erode as
these temporal responses rise.  Yet each next response contains
\(\nu\Delta V_r\), hence two spatial derivatives of the preceding response.
That does not give regularity on its own.  It tells us which spatial
coordinates the equation exposes.

To measure the spatial order exposed by \(\Delta V_r\), let
\(\Lambda=(-\Delta)^{1/2}\) and set
\[
 E_s(t):=\frac12\|\Lambda^su(t)\|_2^2,
 \qquad
 H(t):=E_{1/2}(t),
\]
and retain the complete pressure-coupled nonlinear transfer
\[
 \mathcal P_s(t)
 :=-\operatorname{Re}\left\langle
 \Lambda^su,
 \Lambda^s\bigl((u\!\cdot\nabla)u+\nabla p\bigr)
 \right\rangle_{L^2}.
\]
Pairing the equation at height \(s\) gives
\[
 E_s'=\mathcal P_s-2\nu E_{s+1}.
\]
Differentiate this identity and use it again at each new height.  At every
finite order \(n\),
\[
 H^{(n)}
 =(-2\nu)^nE_{n+1/2}
 +\sum_{j=0}^{n-1}(-2\nu)^j
   \partial_t^{\,n-1-j}\mathcal P_{j+1/2}.
\]
Acceleration, jerk, and the higher temporal responses have now exposed the
Sobolev coordinates
\[
 E_{n+1/2}
 =\frac12\|u\|_{\dot H^{n+1/2}}^2.
\]
If every finite exposed height reaches one terminal value \(u_*\), with the
finite-energy row retained, then for any \(k\) one can choose \(n>k+1\) and
read
\[
 u_*\in L^2\cap\bigcap_{n<\infty}\dot H^{n+1/2}
 =H^\infty
 \hookrightarrow C^\infty.
\]
But exposure is not control.  Differentiation does not bound
\(H^{(n)}\), the transfer terms, or \(E_{n+1/2}\); it only keeps their exact
relation visible.  The singular picture predicts loss of spatial regularity
as temporal order rises, and these identities alone do not prevent it.

Let \([0,T_*)\) be the maximal classical lifespan and suppose
\(T_*<\infty\).  Fix one finite pair \((N,M)\) and one cutoff
\(0<T<T_*\).  For \(0\le n\le N\), keep the whole transfer
\[
 \mathcal B_n
 :=\sum_{a=0}^{n}\binom na(V_a\!\cdot\nabla)V_{n-a}+\nabla Q_n.
\]
Pairing the recurrence for \(V_n\) at height \(M\) gives
\[
\begin{aligned}
 \frac12\frac{d}{dt}\|\Lambda^M V_n\|_2^2
 &+\operatorname{Re}\left\langle
   \Lambda^M V_n,\Lambda^M\mathcal B_n
  \right\rangle_{L^2}\\
 &=\nu\operatorname{Re}\left\langle
   \Lambda^M V_n,\Lambda^M\Delta V_n
  \right\rangle_{L^2}
 =-\nu\|V_n\|_{\dot H^{M+1}}^2.
\end{aligned}
\]
The viscous term has a fixed sign, while every product-rule transfer and the
pressure response remain coupled to it in this finite system.  This identity
shows what must be controlled; it does not supply the whole-terminal bound by
itself.  Theorem~\ref{thm:datum-rectangles} supplies the bound for this maximal
history:
\[
 \sup_{0<T<T_*}\sum_{n=0}^{N}\left(
  \|V_n\|_{L^2(0,T;H^{M+1})}^2
  +\|V_{n+1}\|_{L^2(0,T;H^{M-1})}^2
 \right)<\infty.
\]
Every term in the sum increases monotonically with the cutoff.  Since one
finite bound holds for all \(T<T_*\), the limit \(T\uparrow T_*\) places each
adjacent pair on the full approach to \(T_*\).  With
\(V_{n+1}=\partial_tV_n\), that pair sits in the Hilbert triple
\[
 H^{M+1}\subset H^M\subset H^{M-1}
\]
and gives
\[
 V_n\in C([0,T_*];H^M).
\]
The pressure formula supplies
\(Q_n\in L^1(0,T_*;H^M)\) and
\(\nabla Q_n\in L^1(0,T_*;H^{M-1})\), so the finitely many rows now form one
pressure-complete rectangle
\[
 \mathcal R_{N,M}[u_0;T_*]
 =\bigl((V_n)_{n=0}^{N},(V_{n+1})_{n=0}^{N};(Q_n)_{n=0}^{N}\bigr).
\]
The bound may depend on \(N\) and \(M\).  No constant uniform over every
derivative order is being used.

A larger rectangle retains each older coordinate as the identical
distributional derivative of \(u\) or \(p\).  Restricting from
\((N,M)\) to \((N',M')\) consequently gives
\[
 \rho_{N',M'}^{N,M}\mathcal R_{N,M}[u_0;T_*]
 =\mathcal R_{N',M'}[u_0;T_*].
\]
The restriction maps act between the finite coordinate spaces
\(\mathscr X_{N,M}\) determined by the three rows already displayed.
\[
 \mathscr X_\infty:=\varprojlim_{N,M}\mathscr X_{N,M},
 \qquad
 \mathcal I[u_0]
 :=\bigl(\mathcal R_{N,M}[u_0;T_*]\bigr)_{N,M\in\mathbb N_0}
 \in\mathscr X_\infty,
 \qquad
 u\in\bigcap_{r,m<\infty}H^r(0,T_*;H^m(\mathbb R^3)).
\]
There is no all-order estimate hidden in this notation.  Each requested
coordinate is obtained at a finite stage, and compatibility makes all of
those finite coordinates parts of one history.

Take the zeroth temporal row after this intersection has been formed.  For
every finite \(m\),
\[
 u\in L^2(0,T_*;H^{m+1}),
 \qquad
 u_t\in L^2(0,T_*;H^{m-1}),
 \qquad
 u\in C([0,T_*];H^m).
\]
The endpoint values at different heights are compatible traces of one field.
They produce a common endpoint
\[
 u_*:=u(T_*)\in\bigcap_{m<\infty}H^m(\mathbb R^3)
 =H^\infty(\mathbb R^3)
 \hookrightarrow C^\infty(\mathbb R^3).
\]
At \(T_*\), set \(V_0^*:=u_*\).  The pressure formula and binomial recurrence
then define \(Q_r^*\) and \(V_{r+1}^*\) through every finite order.  Smooth
local existence restarts the flow from \(u_*\).  Uniqueness glues that
continuation to the history already reaching \(T_*\), contradicting finite
maximality.
Hence
\[
 T_*=\infty.
\]

Only after \(\mathcal I[u_0]\) is present do the familiar finite estimates
become available as coordinate projections.  Fix any finite \(T\), a temporal
order \(r\), a spatial derivative order \(k\), and \(2\le p\le\infty\), then
choose
\[
 m>k+\frac32-\frac3p.
\]
The corresponding row gives
\[
 V_r\in C([0,T];H^m)
 \hookrightarrow L^\infty(0,T;W^{k,p}).
\]
For instance, the \(m=3\) coordinate yields
\[
 \int_0^{T}\|\nabla u(t)\|_\infty\,dt
 \le T\|u\|_{L^\infty(0,T;W^{1,\infty})}<\infty.
\]
This estimate appears because the fluid already occupies the compatible
intersection; it is a finite readout of that object, not the source of it.
Its source remains visible in the full dependence
\[
\begin{gathered}
 u_0
 \longrightarrow
 \underbrace{\{(V_r,Q_r)\}_{r\ge0}}_{\text{pressure-complete mixed jet}}
 \longrightarrow
 \underbrace{\left\{
  \begin{gathered}
   V_n\in L^2(0,T_*;H^{M+1}),\\[-1pt]
   V_{n+1}\in L^2(0,T_*;H^{M-1})
  \end{gathered}
 \right\}_{\substack{0\le n\le N\\N,M<\infty}}}_{
 \text{whole-terminal adjacent rows}}
 \\
 \longrightarrow
 \underbrace{\{\mathcal R_{N,M}[u_0;T_*]\}_{N,M<\infty}}_{
 \text{finite compatible rectangles}}
 \longrightarrow
 \underbrace{\mathcal I[u_0]}_{\text{projective intersection}}
 \longrightarrow
 \underbrace{\left(
  u_*\in H^\infty,
  \{(V_r^*,Q_r^*)\}_{r\ge0}
 \right)}_{\substack{\text{common endpoint and pressure-complete}\\
                     \text{endpoint jet}}}
 \\
 \longrightarrow \text{restart and uniqueness}
 \longrightarrow T_*=\infty.
\end{gathered}
\]

\endgroup
